\section{Development Process: Phased Implementation}
\label{sec:process}

\proj{} was built incrementally, one phase at a time, with each phase shipped and
tested before the next began. Two principles held from the start: the original
PDF is \textbf{never modified}, and the headline capability, \textbf{finding text
inside images}, had to be as close to ``100\% findable'' as possible. A third
arrived late and reshaped the whole system: the user must see \textbf{nothing but
Chrome's own viewer}. This section records how the work was planned, how it
actually unfolded, and the reasoning at each step.

\subsection{Planning: The Initial Estimate}
Before writing code, the work was broken into seven phases with rough effort
estimates (Table~\ref{tab:planned}). Two were flagged up front as the likely
time-sinks: the \emph{invisible text-layer overlay} (the coordinate-alignment
problem) and \emph{intercepting Chrome's built-in PDF viewer}, which is sandboxed
and cannot simply be scripted.

\begin{table}[ht]
  \centering
  \caption{Initial seven-phase plan drawn up before implementation began. Phases~5
           and~6 were flagged as the expected time-sinks; Phase~6 was ultimately
           pulled forward and solved first.}
  \label{tab:planned}
  \renewcommand{\arraystretch}{1.2}
  \begin{tabularx}{\textwidth}{@{}clX@{}}
    \toprule
    \textbf{\#} & \textbf{Planned phase} & \textbf{Note} \\
    \midrule
    1 & Extension scaffolding (MV3) & Manifest, service worker, content-script wiring \\
    2 & PDF.js rendering pipeline & Pages to canvas; worker setup in an extension context \\
    3 & Tesseract.js OCR & Plumbing OCR into the async page pipeline \\
    4 & Per-page text detection & \code{getTextContent()} check --- quick \\
    5 & Invisible text-layer overlay & \emph{Hardest:} PDF$\leftrightarrow$canvas$\leftrightarrow$DOM coordinate alignment \\
    6 & Intercept Chrome's PDF viewer & \emph{Hidden complexity:} sandboxed viewer; must redirect \\
    7 & Polish, edge cases, testing & --- \\
    \bottomrule
  \end{tabularx}
\end{table}

\subsection{Execution: How the Build Actually Unfolded}
Execution diverged from the plan in two instructive ways. First, viewer
interception (planned as Phase~6) turned out to be a \emph{prerequisite} for
everything else---you cannot render or overlay anything until you control the
page---so it was solved first, folded into Phase~1. Second, ``OCR accuracy''
was not a single step but the dominant theme of the back half of the project: it
expanded into several phases as the README roadmap and real test feedback
exposed new failure modes. The phases below are the build as it actually
happened.

\medskip
\noindent\textbf{Foundation.}
\begin{description}[leftmargin=1.4em, style=nextline, itemsep=4pt]
  \item[Phase 1 --- Extension skeleton \& PDF interception.] Built the MV3
        manifest, service worker, content script, and an empty viewer, then made
        the service worker redirect \code{.pdf} navigations to that viewer.
        \emph{Why first:} Chrome opens PDFs in its own sandboxed viewer that
        scripts cannot reach, so gaining control of the page is the true
        starting point---and it establishes the read-only-overlay architecture.
        \emph{Revealed:} a viewer that shows nothing yet $\rightarrow$ Phase~2.
  \item[Phase 2 --- PDF.js rendering pipeline.] Bundled PDF.js locally (to satisfy
        the MV3 CSP) and rendered each page to a canvas, switching to a
        \emph{dynamic} redirect rule so the original URL survives as
        \code{?url=}. \emph{Friction:} \code{file://} security origins and the
        \code{getDocument} URL parameter both needed care. \emph{Revealed:} pages
        display but nothing is searchable $\rightarrow$ Phase~3.
  \item[Phase 3 --- Text detection, OCR \& the invisible overlay.] The heart of
        the project: detect native text with \code{getTextContent()}, OCR image
        pages with Tesseract.js in a worker, and inject invisible spans over the
        canvas. \emph{Key pivot:} the first design was \emph{either/or} (use
        native text, else OCR), but a page containing the same word both as native
        text \emph{and} inside an image matched only the native copy---image text
        was lost. The design changed to \emph{both}: always inject native text,
        then OCR images and merge only the non-duplicate words (text-aware
        de-duplication). \emph{Revealed:} OCR re-ran on every load $\rightarrow$
        Phase~4.
  \item[Phase 4 --- IndexedDB caching.] Cache each page's words once, keyed per
        document, with coordinates normalized to $[0,1]$ so the cache survives
        any zoom level. \emph{Why:} OCR is the expensive step, so it should be a
        one-time cost. \emph{Revealed:} language data still implied a CDN
        $\rightarrow$ Phase~5.
  \item[Phase 5 --- Fully offline.] Bundled the LSTM language model and pointed
        Tesseract's \code{langPath} at it, removing the last network dependency.
        \emph{Why:} privacy, reliability, and offline use---no document data ever
        leaves the machine.
\end{description}

\noindent\textbf{Refinement} (working through the README roadmap and test feedback).
\begin{description}[leftmargin=1.4em, style=nextline, itemsep=4pt]
  \item[Phase 6 --- OCR accuracy: correction variants.] Added confidence
        filtering, edge-noise trimming, and digit$\rightarrow$letter correction
        variants (\code{he11o}$\rightarrow$\code{hello}) injected as search-only
        spans. \emph{Thought:} native find is exact-match, so rather than
        fuzzy-matching at query time, widen the \emph{stored} text.
  \item[Phase 7 --- Rotated pages.] Derived each glyph's angle and size from the
        \code{viewport.transform}\,$\times$\,\code{item.transform} matrix and
        rotated the spans to match, so rotation is handled by the math rather
        than special cases (Section~\ref{sec:design}).
  \item[Phase 8 --- Multi-language OCR.] Added an English/Danish selector with a
        per-language cache, so switching languages reprocesses independently and
        non-destructively.
  \item[Phase 9 --- Performance on large PDFs.] Introduced lazy,
        viewport-driven rendering, idle-time background indexing, and
        \emph{parallel} OCR via a worker scheduler. \emph{Thought:} a single
        serial OCR worker was the real bottleneck; the roadmap's ``GPU
        acceleration'' proved not actionable (Tesseract.js has no GPU path and
        already uses the SIMD WASM core), so parallelism was the available win.
  \item[Phase 10 --- Readability \& search UX.] Driven directly by user testing:
        the text layer was made selectable/copyable; a measure-and-scale fix
        aligned each invisible span to its exact width after selection and search
        landed on the wrong characters (``\,Where is\,''$\rightarrow$``\,Where is
        the b\,''); and the native find bar---which flashed the invisible layer
        as oversized doubled text---was replaced by a custom find bar with
        translucent highlights, $\uparrow/\downarrow$ navigation, and a distinct
        active-match style for colourful pages.
  \item[Phase 11 --- OCR image preprocessing.] The final push on the
        ``100\% findable'' goal, targeting the cases still failing: sparse-text
        segmentation recovered \emph{tables}; a high-resolution (3000\,px) OCR
        canvas recovered \emph{small} text; and a percentile contrast stretch
        plus unsharp mask attacked \emph{low-contrast and blurred} text (the
        cache format version was bumped to invalidate stale results).
        \emph{Thought:} decouple the OCR image from the display image and push
        recognition as far as preprocessing allows---while accepting a hard
        ceiling on genuinely low-resolution source images.
\end{description}

\medskip
\noindent\textbf{Production} (preparing for release).
\begin{description}[leftmargin=1.4em, style=nextline, itemsep=4pt]
  \item[Phase 12 --- Activation model: from automatic to on demand.] The biggest
        late change reversed Phase~1's founding decision. Until now the extension
        intercepted \emph{every} PDF navigation with \code{declarativeNetRequest}
        and forced it into the custom viewer. Dogfooding on ordinary text PDFs
        exposed the flaw: Chrome's native PDF viewer is a sandboxed plugin, so
        replacing it threw away its entire toolbar and right-click menu
        (thumbnails, rotate, print, download, Google~Lens, $\dots$)---none of which
        a DOM viewer can reproduce---while for an already-searchable text PDF the
        extension added nothing. \emph{Decision:} keep Chrome's native viewer as
        the default and invoke the OCR viewer \emph{on demand} via a keyboard
        shortcut (\code{Ctrl+Shift+F}) or the toolbar icon, toggling back on a
        second press. \emph{Why:} the extension's value is real only for image and
        scanned PDFs, and those are precisely the documents where the native
        toolbar matters least---so scoping activation to them preserves the native
        experience everywhere else (decision~\ref{dec:ondemand}). \emph{Side
        benefit:} dropping interception also removed the \code{declarativeNetRequest}
        and \code{storage} permissions and the catch-all content script, leaving a
        much smaller, easier-to-review permission footprint.
  \item[Phase 13 --- Imitating the native viewer.] Even on demand, the OCR viewer
        was still a \emph{different} viewer, so the next effort made it look and
        behave like Chrome's: a Chrome-style toolbar with zoom and page
        navigation, a find bar restyled after Chrome's (match counter,
        $\uparrow/\downarrow$, Esc/Enter/Shift+Enter), Chrome's yellow and orange
        highlights, the find bar opening already focused, and a return to the
        native viewer at the same page via \code{\#page=N}. A new search engine
        matched phrases across OCR word boundaries (the old one matched only
        inside a single span), and background indexing stopped freezing in hidden
        tabs once rendering used PDF.js's \code{print} intent, which is not paced
        by \code{requestAnimationFrame}. \emph{Revealed:} the user's verdict was
        that an imitation is not the original---they wanted ``same behaviour as
        that of original \dots\ and same design and view also'', with no
        different view and no change to the PDF's format $\rightarrow$ Phase~14.
  \item[Phase 14 --- Native viewer only: the searchable-scan sandwich.] The
        custom viewer was deleted outright (about 1{,}300 lines). In its place, an
        offscreen document OCRs the PDF and uses pdf-lib to write the recognized
        words into a \emph{copy} as invisible text (\code{3~Tr}), each word sized
        and squeezed so Chrome's highlight covers it (Section~\ref{sec:coords});
        the tab then switches to that copy in Chrome's own viewer. Progress moved
        to a badge on the toolbar icon and the OCR language to the icon's
        right-click menu. \emph{Thought:} the only way to give the user Chrome's
        exact interface is to give Chrome a file it can search itself. This also
        reinterpreted ``never modify the PDF'': the original is still never
        touched, and the copy's visible content was verified identical by pixel
        comparison (0 of 2{,}298{,}825 pixels differed). \emph{Gap:} the pipeline
        was tested in a browser tab, but the final tab switch was not verified in
        Chrome itself.
  \item[Phase 15 --- First use in the user's browser.] Loading the build into the
        user's own Chrome surfaced two problems. OCR crashed inside the offscreen
        document in Tesseract's \emph{relaxed-SIMD} WebAssembly core: Tesseract.js
        chooses that core when a \code{WebAssembly.validate()} probe accepts it,
        but on this Chrome and Apple Silicon combination the probe passed and
        execution trapped. The vendored worker was patched so the probe always
        reports relaxed SIMD as unavailable, falling back to the mature SIMD core.
        Separately, \code{chrome://extensions/shortcuts} showed the Mac shortcut
        (\code{Cmd+Shift+F}) as ``Not set''; that was only explained in
        Phase~16.
  \item[Phase 16 --- Release hardening, verified in real Chrome.] A review before
        publishing to the Chrome Web Store added, for the first time, an
        automated end-to-end test that loads the \emph{packaged} extension into a
        throwaway Chrome profile and drives it through the DevTools
        protocol~\cite{cdp}. It found that the core feature had never worked:
        Chrome silently refuses to navigate a website tab to an extension's
        \code{blob:} URL, so the tab never switched to the searchable copy. Copies
        are now served from the extension's own origin by the service worker
        (decision~\ref{dec:serve}). The same review and tests found and fixed:
        \begin{itemize}[nosep]
          \item OCR collapsing into noise on sparse pages: the contrast stretch's
                dark point fell inside the background, amplifying JPEG noise up
                to $127\times$; the stretch gain is now capped
                (Section~\ref{sec:impl}).
          \item A release zip that omitted the \code{offscreen/} folder, because
                the packaging script still listed the deleted viewer.
          \item The Mac shortcut: Chrome leaves the suggested \code{Cmd+Shift+F}
                unassigned (\code{chrome.commands.getAll()} reports no shortcut),
                so the default is now Control+Shift+F, making the shortcut
                \code{Ctrl+Shift+F} on every platform.
          \item PDFs whose URL lacks ``.pdf'' were ignored; the extension now
                accepts any web or file URL and checks the file's \code{\%PDF-}
                signature.
          \item Danish letters were stripped from word edges (``på'' was stored
                as ``p''), because word trimming only recognized \code{a-z}; it
                now uses Unicode letter classes.
          \item OCR results keyed by URL could be applied to a changed file; the
                cache now keys by a SHA-256 of the content and prunes itself
                (decision~\ref{dec:hash}).
          \item Silent failures: encrypted PDFs, non-PDF pages, and disabled file
                access now show a red badge with an explanation; the success
                badge is set only after the copy loads, because Chrome clears a
                tab's badge when it navigates.
          \item Unbounded resources: pdf.js documents are destroyed after each
                job and idle OCR workers are released
                (decision~\ref{dec:bounded}); unused Tesseract builds were
                removed, shrinking the package from 22.5\,MB to 8.5\,MB.
          \item Store requirements: a manifest description within 132 characters,
                a privacy policy, and listing material.
        \end{itemize}
        \emph{Thought:} unit-level checks had proved the pipeline; only driving
        the real browser proved the product.
\end{description}

\subsection{Plan versus Reality}
The original instinct that the \emph{overlay} would be the hardest part was
correct---coordinate alignment, the native-versus-image de-duplication, and the
measure-and-scale fitting absorbed the most iteration. The estimate was wrong in
its \emph{shape}, though: ``OCR'' looked like one tidy phase but became the
project's centre of gravity, because the headline promise (find \emph{any} text
in \emph{any} image) is exactly the kind of goal that surfaces a long tail of
real-world failure modes---rotation, tables, multiple languages, blur, low
contrast, small type. Each was its own phase, and each was reached by shipping
the previous one and watching where it still fell short.

The plan was also wrong about where the project would end. It assumed a custom
viewer throughout; in the end the custom viewer, its overlay, and its find bar
were all deleted, and the coordinate work moved from the DOM into the PDF itself.
The hardest late problems were not algorithms but \emph{platform behaviour}: what
Chrome lets an extension do with its sandboxed viewer, with blob URLs, with
keyboard shortcuts, and with toolbar badges. None of these showed up in tests of
the pipeline alone; each surfaced only when the real browser was in the loop.
